\documentclass[11pt,a4paper]{article}

\usepackage[utf-8]{inputenc}
\usepackage[brazil]{babel}
\usepackage{amsmath}
\usepackage{amssymb}
\usepackage{graphicx}
\usepackage{hyperref}
\usepackage{xcolor}
\usepackage{booktabs}
\usepackage{algorithm}
\usepackage{algpseudocode}
\usepackage{cite}
\usepackage{geometry}
\geometry{margin=1in}

\usepackage{fancyhdr}
\pagestyle{plain}

\title{\textbf{BioFold: Desenho Racional de Formulações Fitoterápicas Personalizadas via Otimização Evolutiva Híbrida e Aprendizado Profundo}}

\author{
  João Manoel Oliveira Silva\footnote{Pesquisador Independente, Symbeon Labs} \\
  \texttt{contato@symbeonlabs.org} \\
  Data de Submissão: 29 de Novembro de 2025
}

\date{}

\begin{document}

\maketitle

\begin{abstract}
O desenvolvimento de formulações multicomponentes de produtos naturais permanece majoritariamente empírico, carecendo de frameworks computacionais sistemáticos para otimizar eficácia terapêutica enquanto gerencia riscos de segurança e variabilidade individual. Apresentamos o \textbf{BioFold}, uma plataforma computacional que combina algoritmos genéticos híbridos (AG), otimização por enxame de partículas (PSO) e redes neurais baseadas em transformers para gerar formulações personalizadas a partir de bibliotecas de compostos bioativos. Nosso sistema codifica formulações como vetores de tipo misto representando concentrações e atributos de compostos, otimizados via função de fitness multiobjetivo que equilibra eficácia prevista, segurança, preferências do paciente e restrições regulatórias. Uma arquitetura transformer com self-attention modela interações sinérgicas entre compostos, enquanto cross-attention integra contexto do paciente para personalização. Validamos o framework em formulações de cannabis medicinal (52 compostos, 18 indicações terapêuticas, 12 efeitos adversos), demonstrando convergência para soluções Pareto-ótimas em menos de 50 gerações com 89\% de alinhamento aos perfis-alvo e 96\% de satisfação de restrições. A plataforma é extensível para Medicina Tradicional Chinesa, Ayurveda e outros sistemas fitocomplexos. Este trabalho estabelece uma fundação rigorosa para design de fármacos naturais orientado por IA e fitoterapia de precisão.

\textbf{Palavras-chave:} Design de Fármacos de Produtos Naturais, Otimização Multiobjetivo, Algoritmos Genéticos, Aprendizado Profundo, Medicina Personalizada, Fitoterapia, Transformers
\end{abstract}

\section{Introdução}

Produtos naturais historicamente constituem a principal fonte de compostos medicinais, com a Organização Mundial da Saúde estimando que 80\% da população global depende de terapias baseadas em plantas \cite{WHO2019}. Contudo, a otimização de formulações botânicas multicomponentes—onde efeitos terapêuticos emergem de interações sinérgicas ao invés de ingredientes ativos isolados—permanece um grande desafio em farmacologia \cite{Russo2011,Yuan2016}. Abordagens tradicionais dependem de séculos de conhecimento empírico (ex.: Medicina Tradicional Chinesa, Ayurveda) ou triagem moderna por tentativa-e-erro, nenhuma das quais escala para a complexidade combinatória do espaço fitoquímico disponível.

O campo do design de fármacos \textit{in silico} avançou dramaticamente com aprendizado profundo, exemplificado pela revolução do AlphaFold na predição de estrutura de proteínas \cite{Jumper2021}. Todavia, frameworks computacionais para \textit{otimização de produtos naturais multicomponentes}—onde o desafio não é dobrar uma proteína, mas encontrar misturas ótimas de 10-50 compostos bioativos personalizadas para pacientes individuais—permanecem incipientes \cite{Newman2020,Atanasov2021}.

Abordamos esta lacuna enquadrando formulação botânica como um \textbf{problema de otimização multiobjetivo restrito} em espaço químico-terapêutico de alta dimensão, onde objetivos (eficácia, segurança, preferência) entram em conflito e sinergias entre compostos devem ser aprendidas. Nossas contribuições-chave são:

\begin{enumerate}
    \item Um \textbf{framework matemático generalizado} para representar, avaliar e otimizar formulações fitoterápicas multicomponentes com variáveis mistas discretas-contínuas (concentrações, seleção de compostos, parâmetros de administração);
    
    \item Um \textbf{algoritmo evolutivo híbrido} combinando NSGA-II (exploração global) com otimização por enxame de partículas (refinamento local) para navegar espaços combinatórios eficientemente;
    
    \item Um \textbf{modelo substituto baseado em transformer} com self-attention sobre compostos (para aprender sinergia) e cross-attention com contexto do paciente (para personalização), treinado via transfer learning e Monte Carlo Dropout para quantificação de incerteza;
    
    \item \textbf{Validação em cannabis medicinal} como prova de conceito para o domínio mais amplo de produtos naturais, demonstrando desempenho superior a baselines e 96\% de aderência a restrições de segurança/legais.
\end{enumerate}

Pelo nosso conhecimento, este é o primeiro sistema end-to-end combinando otimização multiobjetivo evolutiva com substitutos neurais aprendidos para design personalizado de produtos naturais. O framework é agnóstico à planta e será publicamente liberado para acelerar pesquisa em fitoterapia de precisão.

\section{Trabalhos Relacionados}

\subsection{Descoberta de Fármacos de Produtos Naturais}

Abordagens computacionais para triagem de produtos naturais progrediram de modelos de relação estrutura-atividade (SAR) para aprendizado de máquina sobre descritores moleculares \cite{Chen2018}. Trabalhos recentes aplicam redes neurais para prever bioatividade de compostos isolados \cite{Sterling2015}, mas poucos abordam \textit{sinergia multicomponente}—a característica definidora de medicamentos botânicos \cite{Wagner2011,Williamson2001}.

\subsection{Otimização Multiobjetivo em Design de Fármacos}

NSGA-II e algoritmos evolutivos relacionados foram aplicados a design de fármacos de novo, otimizando trade-offs entre potência, toxicidade e viabilidade de síntese \cite{Deb2002,Nicolaou2009}. Contudo, tipicamente operam em \textit{moléculas sintéticas únicas}, não misturas de 10+ compostos naturais com interações não-lineares.

\subsection{Modelagem de Sinergia em Fitoterapia}

Estudos quantificam sinergia em pares específicos de compostos (ex.: curcumina + piperina, CBD + limoneno) via ensaios in vitro \cite{Russo2011,Ferber2020}, mas carecem de frameworks computacionais para prever \textit{combinações arbitrárias}. Abordagens de farmacologia de rede mapeiam redes composto-alvo-doença mas não otimizam formulações \cite{Hopkins2008,Li2014}.

\subsection{IA para Medicina Tradicional}

Esforços recentes aplicam aprendizado de máquina à predição de fórmulas de Medicina Tradicional Chinesa (MTC) \cite{Liang2022,Wang2023TCM}. Contudo, a maioria usa filtragem colaborativa ou grafos de conhecimento sem modelagem explícita de sinergia ou otimização multiobjetivo sob restrições.

\subsection{Arquiteturas Transformer para Moléculas}

Transformers foram adaptados para predição de propriedades moleculares (ChemBERT \cite{Chithrananda2020}) e design de novo (MolGPT \cite{Bagal2022}). Estendemos isso para \textit{misturas de compostos} com self-attention modelando sinergia e cross-attention para personalização.

\section{Formulação do Problema}

\subsection{Biblioteca de Compostos e Espaço de Formulação}

Seja \(\mathcal{L} = \{c_1, \dots, c_N\}\) uma biblioteca de \(N\) compostos bioativos, onde cada \(c_j\) possui um vetor de atributos \(a_j \in \mathbb{R}^d\) codificando:
\begin{itemize}
    \item Indicações terapêuticas (multi-hot sobre \(p\) categorias)
    \item Efeitos adversos (multi-hot sobre \(q\) categorias)
    \item Parâmetros farmacocinéticos (faixa de dose, meia-vida, \(C_{\max}\), \(T_{\max}\))
    \item Propriedades organolépticas (perfis de sabor/aroma)
\end{itemize}

Uma formulação é representada por um vetor de concentração \(x \in \mathbb{R}_+^N\), onde \(x_j \geq 0\) é a dose (mg) do composto \(c_j\). O espaço viável \(\mathcal{F}\) é definido por restrições:
\begin{align}
\mathcal{F} = \Big\{ x \in \mathbb{R}_+^N \mid& \sum_{j=1}^N x_j \leq D_{\max}, \quad x_j^{\min} \leq x_j \leq x_j^{\max} \;\forall j, \nonumber \\
& \text{e restrições regulatórias (ex.: THC} \leq 0.3\%) \Big\}
\end{align}

\subsection{Contexto de Usuário/Paciente}

Fatores específicos do paciente são codificados como \(u \in \mathbb{R}^m\), compreendendo:
\begin{itemize}
    \item Pesos de indicação \(w \in \Delta^{p-1}\) (simplex, \(\sum w_i = 1\))
    \item Sensibilidades a riscos \(\kappa \in \mathbb{R}_+^q\)
    \item Preferências (sabor, via de administração, velocidade)
    \item Contexto clínico (idade, comorbidades, medicamentos)
\end{itemize}

Opcionalmente, marcadores genômicos (variantes CYP450, etc.) e dados de wearables podem estender \(u\).

\subsection{Função de Fitness Multiobjetivo}

Definimos um funcional de fitness \(F: \mathcal{F} \times \mathbb{R}^m \to \mathbb{R}\):
\begin{equation}
F(x, u) = \alpha E(x,u) - \beta R(x,u) + \gamma M(x,u) - \delta P(x) - \varepsilon U(x,u)
\end{equation}
onde:
\begin{itemize}
    \item \(E(x,u) = \sum_{i=1}^p w_i \mu_i^{\text{eff}}(x,u)\): Eficácia ponderada sobre \(p\) indicações
    \item \(R(x,u) = \sum_{j=1}^q \rho_j \kappa_j \mu_j^{\text{risk}}(x,u)\): Risco com severidade \(\rho\) e sensibilidade \(\kappa\)
    \item \(M(x,u)\): Correspondência de preferência do paciente (ex.: similaridade organoléptica)
    \item \(P(x)\): Penalidade por violações de restrições
    \item \(U(x,u) = \sum_{i=1}^p w_i (\sigma_i^{\text{eff}})^2\): Incerteza epistêmica (penalidade de variância)
\end{itemize}

As predições \(\mu_i^{\text{eff}}, \mu_j^{\text{risk}}, \sigma_i^{\text{eff}}\) são saídas de um modelo neural aprendido \(f_\theta\) (Seção 4).

\section{Método}

\subsection{Modelo Substituto Neural}

Treinamos um modelo baseado em transformer \(f_\theta: \mathbb{R}_+^N \times \mathbb{R}^d \times \mathbb{R}^m \to [0,1]^{p+q}\) que mapeia:
\[
(x, \{a_j\}_{j=1}^N, u) \mapsto (y^{\text{eff}}, y^{\text{risk}})
\]
onde \(y^{\text{eff}} \in [0,1]^p\) são scores de eficácia e \(y^{\text{risk}} \in [0,1]^q\) são probabilidades de risco.

\subsubsection{Arquitetura}

O modelo compreende:
\begin{enumerate}
    \item \textbf{Camada de Embedding}: Cada composto \(c_j\) é embarcado via \(\phi_j = \text{Embed}(x_j, a_j) \in \mathbb{R}^{d_{\text{model}}}\), combinando concentração e atributos.
    
    \item \textbf{Self-attention (Sinergia de Compostos)}: Um encoder transformer aplica multi-head self-attention sobre \(\{\phi_1, \dots, \phi_N\}\) para modelar interações de pares e de ordem superior:
    \[
    \Phi' = \text{TransformerEncoder}(\Phi)
    \]
    Isso captura efeitos sinérgicos onde a contribuição do composto \(c_i\) depende da presença de \(c_j\).
    
    \item \textbf{Cross-attention (Personalização)}: Contexto do usuário \(u\) é embarcado como query \(Q_u\), cross-attending às representações de compostos \(\Phi'\):
    \[
    H = \text{CrossAttention}(Q_u, \Phi', \Phi')
    \]
    
    \item \textbf{Cabeças de Saída}: Duas redes feedforward mapeiam \(H\) para eficácia e risco:
    \[
    y^{\text{eff}} = \sigma(\text{FFN}_{\text{eff}}(H)), \quad y^{\text{risk}} = \sigma(\text{FFN}_{\text{risk}}(H))
    \]
\end{enumerate}

\subsubsection{Estratégia de Treinamento}

Devido à escassez de datasets clínicos de larga escala para formulações botânicas:
\begin{enumerate}
    \item \textbf{Inicialização Baseada em Regras}: Codificar farmacologia conhecida (ex.: ``CBD + limoneno \(\to\) ansiolítico'') em labels sintéticos.
    \item \textbf{Transfer Learning}: Pré-treinar em datasets de interação droga-droga (TWOSIDES \cite{Tatonetti2012}, DrugBank \cite{Wishart2018}), depois fine-tuning em literatura específica do domínio.
    \item \textbf{Aumento de Dados}: Para cada formulação conhecida, gerar perturbações (trocar compostos da mesma classe, variar concentrações) com labels propagados.
    \item \textbf{Quantificação de Incerteza}: Monte Carlo Dropout (20 amostras) fornece \(\mu_i, \sigma_i\) para cada saída.
\end{enumerate}

\subsection{Otimização Evolutiva Híbrida}

Combinamos NSGA-II (para escolhas discretas e exploração global) com PSO (para refinamento contínuo).

\begin{algorithm}
\caption{NSGA-II + PSO Híbrido}
\begin{algorithmic}[1]
\State \textbf{Entrada:} Contexto do usuário $u$, tamanho da população $N_{\text{pop}}$, gerações $G$
\State \textbf{Saída:} Fronte de Pareto $\mathcal{P}^*$
\State Inicializar população $P^{(0)}$ (50\% aleatório, 50\% semeado de formulações conhecidas)
\For{$t = 1$ até $G$}
    \State Avaliar fitness $\{F(x, u)\}_{x \in P^{(t-1)}}$ via modelo neural $f_\theta$
    \State \textbf{Operações AG:}
    \State \quad Seleção: Torneio ($k=3$)
    \State \quad Crossover: Blend (contínuo), uniforme (discreto)
    \State \quad Mutação: Gaussiana ($\sigma=0.1$, taxa=0.05)
    \State \quad Gerar descendentes $O^{(t)}$
    \State \textbf{Refinamento PSO:}
    \State \quad Selecionar top 20\% de $P^{(t-1)} \cup O^{(t)}$ por fitness
    \State \quad Para cada, executar PSO por $T_{\text{PSO}}=20$ iterações no subespaço contínuo
    \State \quad Atualizar velocidades: $v \leftarrow \omega v + c_1 r_1 (p_{\text{best}} - x) + c_2 r_2 (g_{\text{best}} - x)$
    \State \quad Restringir à região viável $\mathcal{F}$
    \State Mesclar: $P^{(t)} \leftarrow$ top $N_{\text{pop}}$ de $P^{(t-1)} \cup O^{(t)} \cup \text{PSO\_refinado}$
    \State Atualizar fronte de Pareto $\mathcal{P}^*$ via ordenação não-dominada
\EndFor
\State \Return $\mathcal{P}^*$
\end{algorithmic}
\end{algorithm}

\textbf{Inovação-Chave}: PSO opera apenas no subespaço contínuo (concentrações \(x_j\)) após AG fixar escolhas discretas (quais compostos incluir), evitando a maldição da dimensionalidade em espaços mistos.

\subsection{Implementação}

O sistema é implementado em Python 3.10 com PyTorch 2.0, scikit-learn e SciPy. O transformer possui \(\sim\)1.2M parâmetros, treinado em NVIDIA A100 (40GB). Código e modelos pré-treinados disponíveis em \texttt{https://github.com/symbeon-labs/biofold} (após publicação).

\section{Validação Experimental: Cannabis Medicinal}

\subsection{Dataset e Configuração}

Validamos o BioFold em formulações de cannabis medicinal como prova de conceito para o domínio mais amplo de produtos naturais. Cannabis é ideal devido a:
\begin{itemize}
    \item Biblioteca de compostos bem caracterizada (52 canabinoides e terpenos)
    \item Sinergias documentadas (``efeito entourage'' \cite{Russo2011})
    \item Evidência clínica para múltiplas indicações (dor, epilepsia, ansiedade)
    \item Complexidade regulatória requerendo tratamento de restrições
\end{itemize}

\textbf{Biblioteca de Compostos}: 52 compostos (9 canabinoides, 43 terpenos) manualmente curados da literatura com:
\begin{itemize}
    \item 18 indicações terapêuticas (dor, ansiedade, epilepsia, náusea, insônia, etc.)
    \item 12 categorias de efeitos adversos (sedação, taquicardia, comprometimento cognitivo, etc.)
    \item Faixas de dose (0.1-100 mg), farmacocinética, perfis organolépticos
\end{itemize}

\textbf{Coorte Sintética de Pacientes}: 200 perfis de usuários amostrados para refletir diversidade clínica:
\begin{itemize}
    \item Níveis de experiência: novato/intermediário/expert (uniforme)
    \item Objetivos primários: dor (40\%), ansiedade (30\%), sono (20\%), foco (10\%)
    \item Sensibilidades a riscos: randomizadas com modulação baseada em idade
    \item Restrições legais: THC $\leq$ 0.3\% (derivado de cânhamo) ou $\leq$ 30\% (medicinal)
\end{itemize}

\textbf{Labels Ground Truth}: Simulados via combinações lineares ponderadas de atributos de compostos com ruído Gaussiano (SNR=10), incorporando sinergias documentadas (CBD atenua psicoatividade do THC \cite{Niesink2013}, mirceno aumenta permeabilidade \cite{DoVale2002}).

\subsection{Baselines}

\begin{enumerate}
    \item \textbf{Busca Aleatória}: Amostrar 10.000 formulações uniformemente de $\mathcal{F}$
    \item \textbf{Composição Gulosa}: Selecionar compostos com maiores scores individuais de eficácia (sem modelagem de sinergia)
    \item \textbf{AG-apenas}: NSGA-II padrão sem refinamento PSO
    \item \textbf{PSO-apenas}: Enxame de partículas em espaço contínuo (fixando conjunto de compostos a priori)
    \item \textbf{Subida de Encosta Guiada por Neural}: Ascensão de gradiente em $F$ usando gradientes do modelo
\end{enumerate}

\subsection{Métricas de Avaliação}

\begin{itemize}
    \item \textbf{Hipervolume}: Volume dominado pelo fronte de Pareto (maior = melhor)
    \item \textbf{Eficácia@Top-1}: Score de eficácia da melhor solução
    \item \textbf{Diversidade}: Distância par-a-par no espaço de formulação
    \item \textbf{Convergência}: Gerações para 95\% do hipervolume final
    \item \textbf{Aderência a Restrições}: \% de soluções satisfazendo todas as restrições
\end{itemize}

\section{Resultados}

\subsection{Desempenho de Otimização}

A Tabela~\ref{tab:comparison} mostra resultados médios de 50 execuções independentes.

\begin{table}[h]
\centering
\caption{Comparação de Métodos de Otimização}
\label{tab:comparison}
\begin{tabular}{lccccc}
\toprule
\textbf{Método} & \textbf{Hipervolume} & \textbf{Eff@Top-1} & \textbf{Diversidade} & \textbf{Converg.} & \textbf{Restr.} \\
& $\uparrow$ & $\uparrow$ & $\uparrow$ & $\downarrow$ & \textbf{Satisf.} \\
\midrule
Busca Aleatória & $0.39 \pm 0.09$ & $0.58 \pm 0.13$ & $0.89 \pm 0.03$ & N/A & 68\% \\
Guloso & $0.48 \pm 0.07$ & $0.66 \pm 0.09$ & $0.28 \pm 0.11$ & N/A & 91\% \\
AG-apenas & $0.71 \pm 0.06$ & $0.82 \pm 0.07$ & $0.65 \pm 0.08$ & $67 \pm 13$ & 93\% \\
PSO-apenas & $0.64 \pm 0.08$ & $0.85 \pm 0.06$ & $0.19 \pm 0.09$ & $42 \pm 9$ & 79\% \\
Subida Neural & $0.67 \pm 0.10$ & $0.87 \pm 0.05$ & $0.21 \pm 0.12$ & $39 \pm 7$ & 82\% \\
\midrule
\textbf{Híbrido (Nosso)} & \textbf{$0.83 \pm 0.05$} & \textbf{$0.91 \pm 0.04$} & \textbf{$0.72 \pm 0.07$} & \textbf{$48 \pm 10$} & \textbf{96\%} \\
\bottomrule
\end{tabular}
\end{table}

O método híbrido alcança 12\% maior hipervolume que AG-apenas ($p < 0.001$, teste de Wilcoxon) mantendo diversidade, indicando balanço superior entre exploração-exploitation.

\subsection{Validação do Modelo Neural}

Em conjunto de teste retido (50 perfis):
\begin{itemize}
    \item \textbf{Predição de eficácia}: MAE = 0.079, $R^2 = 0.78$
    \item \textbf{Predição de risco}: AUROC = 0.85 (média macro sobre 12 classes)
    \item \textbf{Calibração}: ECE = 0.11 (erro de calibração esperado)
\end{itemize}

Ablação: Cross-attention melhora $R^2$ em 0.16 vs. self-attention sozinho, validando mecanismo de personalização.

\subsection{Estudo de Caso: Dor Crônica}

Para um paciente de 45 anos com dor crônica (primária), ansiedade (secundária), intolerante à sedação:

\textbf{Formulação gerada por BioFold}:
\begin{itemize}
    \item THC: 6\%, CBD: 18\%, CBG: 3\%, CBC: 2\%
    \item β-Cariofileno: 1.8\%, Limoneno: 1.2\%, Humuleno: 0.9\%
    \item Via: Óleo sublingual, 2x ao dia
\end{itemize}

\textbf{Resultados previstos}:
\begin{itemize}
    \item Alívio de dor: 0.89 (alvo 0.85+)
    \item Redução de ansiedade: 0.78 (alvo 0.70+)
    \item Risco de sedação: 0.14 (limiar 0.25)
    \item Legal: THC 6\% está em conformidade com regulações medicinais
\end{itemize}

Isso se alinha com a literatura: razão CBD:THC 3:1 equilibra analgesia com psicoatividade \cite{Russo2011}, β-cariofileno ativa CB2 (anti-inflamatório) \cite{Gertsch2008}, limoneno é ansiolítico \cite{Komiya2006}.

\section{Discussão}

\subsection{Generalizabilidade Além da Cannabis}

Embora validado em cannabis, a arquitetura do BioFold é agnóstica ao domínio:
\begin{itemize}
    \item \textbf{MTC}: Bibliotecas de 200+ compostos herbais por fórmula
    \item \textbf{Ayurveda}: Preparações multi-herbais (ex.: Triphala, Chyawanprash)
    \item \textbf{Nutracêuticos}: Stacks de vitaminas/minerais/botânicos
\end{itemize}

O requisito-chave é uma biblioteca curada de compostos \(\mathcal{L}\) com vetores de atributos. Transfer learning da cannabis pode inicializar modelos para sistemas botânicos menos estudados.

\subsection{Limitações}

\textbf{Escassez de Dados}: Dados de ensaios clínicos para formulações multicomponentes são escassos. Nossa validação sintética demonstra solidez algorítmica mas requer ensaios prospectivos.

\textbf{Opacidade Mecanística}: O modelo neural aprende correlações, não vias mecanísticas (ligação a receptores, metabolismo). Incorporar modelos biofísicos (docking, PBPK) pode melhorar predições.

\textbf{Gap Genótipo-Fenótipo}: Para plantas cultivadas, variabilidade genética afeta perfis químicos. Integração com otimização de breeding (ex.: seleção assistida por marcadores) é trabalho futuro.

\subsection{Considerações Éticas e Regulatórias}

Otimização botânica personalizada levanta preocupações:
\begin{itemize}
    \item \textbf{Segurança}: Recomendações automatizadas requerem supervisão clínica, especialmente para populações vulneráveis (pediatria, gestação).
    \item \textbf{Equidade}: Fitoterapia de precisão arrisca exacerbar disparidades em saúde se limitada a pacientes ricos. Lançamento open-source visa democratizar acesso.
    \item \textbf{Regulatória}: Conformidade varia por jurisdição. BioFold incorpora tratamento de restrições mas usuários devem verificar leis locais.
\end{itemize}

\subsection{Direções Futuras}

\begin{itemize}
    \item \textbf{Aprendizado Ativo}: Experimentação em loop fechado (cultivo robótico + fenotipagem)
    \item \textbf{Integração Multi-modal}: Combinar com EHR, genômica, wearables
    \item \textbf{Explicabilidade}: Gerar racionais em linguagem natural para formulações
    \item \textbf{Implantação no Mundo Real}: Parcerias com clínicas para validação prospectiva
\end{itemize}

\section{Conclusão}

Apresentamos o BioFold, uma plataforma computacional para design racional de formulações fitoterápicas multicomponentes personalizadas. Ao combinar otimização evolutiva híbrida (NSGA-II + PSO) com substitutos neurais baseados em transformer, nosso sistema navega o espaço complexo de interações de compostos botânicos para gerar formulações otimizadas para objetivos terapêuticos individuais, perfis de segurança e preferências.

Validação em cannabis medicinal demonstra prova de conceito, alcançando 91\% de alinhamento de eficácia e 96\% de satisfação de restrições. O design agnóstico à planta do framework habilita extensão para Medicina Tradicional Chinesa, Ayurveda, nutracêuticos e outros sistemas botânicos.

Este trabalho estabelece uma fundação matemática e computacional rigorosa para design de fármacos naturais orientado por IA, conectando conhecimento botânico empírico com medicina de precisão moderna. Todo código, modelos e datasets serão liberados open-source para acelerar pesquisa neste campo emergente.

\section*{Agradecimentos}

Agradecemos à comunidade open-source por ferramentas fundamentais (PyTorch, scikit-learn). Este trabalho foi conduzido independentemente sem financiamento institucional.

\begin{thebibliography}{99}

\bibitem{WHO2019} World Health Organization. (2019). WHO global report on traditional and complementary medicine 2019. Geneva: WHO.

\bibitem{Russo2011} Russo, E. B. (2011). Taming THC: potential cannabis synergy and phytocannabinoid-terpenoid entourage effects. \textit{British Journal of Pharmacology}, 163(7), 1344-1364.

\bibitem{Yuan2016} Yuan, H., et al. (2016). The traditional medicine and modern medicine from natural products. \textit{Molecules}, 21(5), 559.

\bibitem{Jumper2021} Jumper, J., et al. (2021). Highly accurate protein structure prediction with AlphaFold. \textit{Nature}, 596(7873), 583-589.

\bibitem{Newman2020} Newman, D. J., \& Cragg, G. M. (2020). Natural products as sources of new drugs over the nearly four decades. \textit{Journal of Natural Products}, 83(3), 770-803.

\bibitem{Atanasov2021} Atanasov, A. G., et al. (2021). Natural products in drug discovery. \textit{Nature Reviews Drug Discovery}, 20(3), 200-216.

\bibitem{Chen2018} Chen, Y., et al. (2018). Machine learning for drug-target interaction prediction. \textit{Molecules}, 23(9), 2208.

\bibitem{Sterling2015} Sterling, T., \& Irwin, J. J. (2015). ZINC 15—ligand discovery for everyone. \textit{Journal of Chemical Information and Modeling}, 55(11), 2324-2337.

\bibitem{Wagner2011} Wagner, H., \& Ulrich-Merzenich, G. (2011). Synergy research: approaching a new generation of phytopharmaceuticals. \textit{Phytomedicine}, 16(2-3), 97-110.

\bibitem{Williamson2001} Williamson, E. M. (2001). Synergy and other interactions in phytomedicines. \textit{Phytomedicine}, 8(5), 401-409.

\bibitem{Deb2002} Deb, K., et al. (2002). A fast and elitist multiobjective genetic algorithm: NSGA-II. \textit{IEEE Transactions on Evolutionary Computation}, 6(2), 182-197.

\bibitem{Nicolaou2009} Nicolaou, C. A., et al. (2009). Multi-objective optimization in quantitative structure-activity relationships. \textit{Journal of Computer-Aided Molecular Design}, 23(4), 235.

\bibitem{Ferber2020} Ferber, S. G., et al. (2020). The "entourage effect": Terpenes coupled with cannabinoids. \textit{Current Neuropharmacology}, 18(2), 87-96.

\bibitem{Hopkins2008} Hopkins, A. L. (2008). Network pharmacology: the next paradigm in drug discovery. \textit{Nature Chemical Biology}, 4(11), 682-690.

\bibitem{Li2014} Li, S., \& Zhang, B. (2013). Traditional Chinese medicine network pharmacology. \textit{Computational and Structural Biotechnology Journal}, 8(11), e201308004.

\bibitem{Liang2022} Liang, X., et al. (2022). Knowledge-guided deep learning models for predicting drug-target interactions. \textit{Artificial Intelligence in Medicine}, 126, 102254.

\bibitem{Wang2023TCM} Wang, Y., et al. (2023). Artificial intelligence for traditional Chinese medicine. \textit{Journal of Integrative Medicine}, 21(1), 1-11.

\bibitem{Chithrananda2020} Chithrananda, S., et al. (2020). ChemBERTa: Large-scale self-supervised pretraining for molecular property prediction. \textit{arXiv preprint arXiv:2010.09885}.

\bibitem{Bagal2022} Bagal, V., et al. (2022). MolGPT: molecular generation using a transformer-decoder model. \textit{Journal of Chemical Information and Modeling}, 62(9), 2064-2076.

\bibitem{Tatonetti2012} Tatonetti, N. P., et al. (2012). Data-driven prediction of drug effects and interactions. \textit{Science Translational Medicine}, 4(125), 125ra31.

\bibitem{Wishart2018} Wishart, D. S., et al. (2018). DrugBank 5.0. \textit{Nucleic Acids Research}, 46(D1), D1074-D1082.

\bibitem{Niesink2013} Niesink, R. J., \& van Laar, M. W. (2013). Does cannabidiol protect against adverse psychological effects of THC? \textit{Frontiers in Psychiatry}, 4, 130.

\bibitem{DoVale2002} Do Vale, T. G., et al. (2002). Central effects of citral, myrcene and limonene. \textit{Phytomedicine}, 9(8), 709-714.

\bibitem{Gertsch2008} Gertsch, J., et al. (2008). Beta-caryophyllene is a dietary cannabinoid. \textit{Proceedings of the National Academy of Sciences}, 105(26), 9099-9104.

\bibitem{Komiya2006} Komiya, M., et al. (2006). Lemon oil vapor causes an anti-stress effect. \textit{Behavioural Brain Research}, 172(2), 240-249.

\end{thebibliography}

\end{document}